% Adaptacja AF-12 i AM I-07; oryginalne etykiety zachowane.
\chapter{Przestrzenie unormowane i przestrzenie Banacha}
\label{ma2:ch:normed}
\providecommand{\K}{\mathbb K}
Niech $\K\in\{\R,\C\}$. Zakładamy znajomość elementarnej
algebry liniowej, ciągów i zupełności metrycznej z tomu I.
\section{Normy i operatory}
\begin{definicja}[Przestrzeń unormowana i przestrzeń Banacha]
\label{def:banach12}
Przestrzenią unormowaną nazywamy parę $(E,\|\cdot\|)$, gdzie $E$
jest przestrzenią wektorową nad $\K$, a norma spełnia
$\|x\|=0\iff x=0$, $\|\alpha x\|=|\alpha|\|x\|$ oraz
$\|x+y\|\le\|x\|+\|y\|$. Przestrzeń unormowaną nazywamy
\emph{przestrzenią Banacha}, jeżeli jest zupełna względem metryki
$d(x,y)=\|x-y\|$.
\end{definicja}

\begin{stwierdzenie}[Norma operatora i przestrzeń sprzężona]
\label{prop:operator-norm12}
Niech $E$ będzie przestrzenią unormowaną i $F$ przestrzenią Banacha.
Wtedy przestrzeń $\mathcal L(E,F)$ wszystkich ciągłych operatorów
liniowych $E\to F$, z normą
$\|T\|=\sup_{\|x\|_E\le1}\|Tx\|_F$, jest przestrzenią Banacha.
W szczególności $E'=\mathcal L(E,\K)$ jest zawsze przestrzenią
Banacha, nawet gdy $E$ nie jest zupełna.
\end{stwierdzenie}
\begin{proof}
Ciągłość liniowego $T$ w zerze jest równoważna istnieniu $C\ge0$
takiego, że $\|Tx\|_F\le C\|x\|_E$: dla implikacji nietrywialnej
wystarczy przeskalować niezerowe $x$ do kuli, na której $T$ jest
ograniczony. Norma operatorowa jest zatem skończona i spełnia
aksjomaty normy.

Jeśli $(T_n)$ jest ciągiem Cauchy'ego w normie operatorowej,
to dla każdego $x\in E$ ciąg $(T_nx)$ jest Cauchy'ego w $F$, ponieważ
$\|T_nx-T_mx\|\le\|T_n-T_m\|\|x\|$.
Niech $Tx=\lim_nT_nx$. Granice zachowują liniowość. Ustalmy
$\varepsilon>0$ i wybierzmy $N$ tak, że
$\|T_n-T_m\|<\varepsilon$ dla $m,n\ge N$.
Po przejściu do granicy $m\to\infty$ otrzymujemy
$\|(T_n-T)x\|\le\varepsilon\|x\|$ dla $n\ge N$.
Stąd $T$ jest ciągły i $\|T_n-T\|\to0$.
\end{proof}

\section{Przestrzenie ciągowe}

\begin{definicja}[Przestrzenie $\ell^p$, $c_0$ i $\ell^\infty$]
\label{def:sequence-spaces12}
Dla $1\le p<\infty$ określamy
$\ell^p=\{x=(x_n)_{n\ge1}\in\K^{\N}:\sum_{n\ge1}|x_n|^p<\infty\}$
i $\|x\|_p=(\sum_{n\ge1}|x_n|^p)^{1/p}$.
Ponadto $\ell^\infty$ jest przestrzenią ciągów ograniczonych
z normą $\|x\|_\infty=\sup_{n\ge1}|x_n|$, a $c_0$ oznacza
podprzestrzeń ciągów zbieżnych do zera, z tą samą normą.
Przez $c_{00}$ oznaczamy przestrzeń ciągów o skończonym nośniku.
\end{definicja}

\begin{lemat}[Nierówność H\"oldera dla ciągów]
\label{lem:holder-seq12}
Dla $1<p<\infty$, $q=p/(p-1)$, $x\in\ell^p$ i $y\in\ell^q$
zachodzi
$\sum_{n\ge1}|x_ny_n|\le\|x\|_p\|y\|_q$.
Dla $p=1$, $q=\infty$, ta sama nierówność jest prawdziwa.
\end{lemat}
\begin{proof}
Dla $a,b\ge0$ nierówność Younga
$ab\le a^p/p+b^q/q$ wynika z wypukłości funkcji
$t\mapsto e^t$ (równoważnie, z minimalizacji
$t^p/p-tb$ względem $t\ge0$). Jeżeli obie normy są niezerowe,
stosujemy ją do $a=|x_n|/\|x\|_p$ i
$b=|y_n|/\|y\|_q$, sumujemy po $n\le N$,
a następnie przechodzimy do granicy $N\to\infty$.
Jeśli jedna norma jest zerowa, wszystkie współrzędne tego ciągu są zerowe, więc obie strony nierówności są zerowe. Dla $p=1$
oszacowujemy bezpośrednio każdy wyraz przez
$|x_n|\|y\|_\infty$.
\end{proof}

\begin{twierdzenie}[Zupełność klasycznych przestrzeni ciągowych]
\label{thm:seq-banach12}
Dla każdego $1\le p\le\infty$ przestrzeń $\ell^p$ jest Banacha.
Przestrzeń $c_0$ jest domkniętą podprzestrzenią $\ell^\infty$,
a zatem także przestrzenią Banacha. Ponadto $c_{00}$ jest gęste
w $\ell^p$ dla $1\le p<\infty$ i w $c_0$, lecz nie jest gęste
w $\ell^\infty$.
\end{twierdzenie}
\begin{proof}
Dla $p=\infty$ norma supremum jest normą: nieujemność i
rozdzielanie punktów wynikają z modułu, jednorodność z
$|\alpha x_j|=|\alpha||x_j|$, a nierówność trójkąta z
$|x_j+y_j|\le\|x\|_\infty+\|y\|_\infty$ dla każdego $j$.
Dla $p<\infty$ jednorodność i rozdzielanie punktów wynikają
bezpośrednio z definicji. Dla $p=1$ nierówność trójkąta
otrzymujemy przez zsumowanie nierówności dla modułu.
Dla $p>1$ wypukłość $t\mapsto t^p$ na $[0,\infty)$ daje
$|x_j+y_j|^p\le2^{p-1}(|x_j|^p+|y_j|^p)$, więc $z=x+y\in\ell^p$.
Niech $q=p/(p-1)$ i $w_j=|z_j|^{p-1}$.
Ponieważ $(p-1)q=p$, mamy $w\in\ell^q$ i
$\|w\|_q=\|z\|_p^{p-1}$. Z lematu \ref{lem:holder-seq12},
\[
 \|z\|_p^p\le\sum_j|x_j|w_j+\sum_j|y_j|w_j
 \le(\|x\|_p+\|y\|_p)\|z\|_p^{p-1}.
\]
Gdy $\|z\|_p=0$, nierówność trójkąta wynika z nieujemności
prawej strony; inaczej dzielimy przez $\|z\|_p^{p-1}>0$.
Wykazaliśmy również domknięcie na dodawanie; domknięcie na
mnożenie przez skalar wynika z jednorodności.
Niech $(x^{(n)})$ będzie ciągiem Cauchy'ego w $\ell^p$.
Dla każdego $j$ zachodzi
$|x_j^{(n)}-x_j^{(m)}|\le\|x^{(n)}-x^{(m)}\|_p$.
Z zupełności $\K$ istnieje $x_j=\lim_m x_j^{(m)}$.
Ustalmy $\varepsilon>0$ i $N$ tak, że norma różnicy jest
mniejsza niż $\varepsilon/2$ dla $m,n\ge N$.
Dla $p<\infty$ i dowolnego $J\ge1$ mamy
\[
 \sum_{j=1}^J|x_j^{(n)}-x_j^{(m)}|^p\le(\varepsilon/2)^p.
\]
Przechodząc do granicy $m\to\infty$ w skończonej sumie,
a następnie do granicy niemalejących sum po $J$, dostajemy
$x^{(n)}-x\in\ell^p$ oraz
$\|x^{(n)}-x\|_p\le\varepsilon/2$ dla $n\ge N$.
Dla $p=\infty$ przechodzimy do granicy w każdej nierówności
$|x_j^{(n)}-x_j^{(m)}|\le\varepsilon/2$, po czym bierzemy
supremum po $j$; otrzymujemy identyczne oszacowanie.
W obu przypadkach $x=x^{(N)}-(x^{(N)}-x)\in\ell^p$,
a oszacowanie daje zbieżność w normie.
Jeśli $x^{(n)}\in c_0$ i $x^{(n)}\to x$ w normie supremum,
wybieramy $n$ z $\|x-x^{(n)}\|_\infty<\varepsilon/2$
oraz $J$ z $|x_j^{(n)}|<\varepsilon/2$ dla $j\ge J$.
Wówczas $|x_j|<\varepsilon$ dla $j\ge J$, więc $x\in c_0$.
Z rachunku granic skalarnych $c_0$ jest podprzestrzenią liniową.
Jako domknięta podprzestrzeń zupełnej przestrzeni jest zupełna.
Obcięcie $x$ po $J$ współrzędnych daje $x^{[J]}\in c_{00}$.
Dla $p<\infty$ norma różnicy wynosi
$(\sum_{j>J}|x_j|^p)^{1/p}\to0$, a dla $x\in c_0$
wynosi $\sup_{j>J}|x_j|\to0$. Dla $u=(1,1,\ldots)$
i każdego $v\in c_{00}$ pewna współrzędna $v_j$ jest zerowa,
więc $\|u-v\|_\infty\ge1$, co wyklucza gęstość w $\ell^\infty$.
\end{proof}



\section{Metryka, równoważność norm i zupełność}
\begin{twierdzenie}
\label{ma2:thm:metric}
Norma wyznacza metrykę $d(x,y)=\|x-y\|$. Norma jest ciągła,
a dodawanie i mnożenie przez skalar są ciągłe w topologiach
produktowych.
\end{twierdzenie}
\begin{proof}
Z $0=\|x+(-x)\|\le2\|x\|$ wynika nieujemność normy.
Rozdzielanie punktów i symetria wynikają z aksjomatów normy,
w tym jednorodności dla skalara $-1$. Nierówność trójkąta
dla metryki wynika z $x-z=(x-y)+(y-z)$.
Z $\|x\|\le\|x-y\|+\|y\|$ oraz nierówności z zamienionymi
$x,y$ otrzymujemy $|\|x\|-\|y\||\le\|x-y\|$.
Jeśli $x_n\to x$, $y_n\to y$, to
$\|x_n+y_n-x-y\|\le\|x_n-x\|+\|y_n-y\|\to0$.
Jeżeli $\alpha_n\to\alpha$, ciąg $(|\alpha_n|)$ jest
ograniczony przez $M$ i
$\|\alpha_n x_n-\alpha x\|\le M\|x_n-x\|+|\alpha_n-\alpha|\|x\|\to0$.
Kryterium ciągowe w przestrzeniach metrycznych daje ciągłość działań.
\end{proof}
\begin{twierdzenie}[Równoważność norm]
\label{ma2:thm:equiv}
Jeżeli dla norm $p,q$ na $E$ istnieją $a,b>0$ z
$a p(x)\le q(x)\le b p(x)$ dla każdego $x\in E$, to normy
wyznaczają te same topologie, zbieżności i ciągi Cauchy'ego.
Zupełność jednej jest równoważna zupełności drugiej.
\end{twierdzenie}
\begin{proof}
Inkluzje $B_p(x,r/b)\subseteq B_q(x,r)$ i
$B_q(x,ar)\subseteq B_p(x,r)$ dają równość topologii.
Stosując oszacowania do $x_n-x$ oraz $x_n-x_m$,
otrzymujemy równoważność odpowiednio warunków zbieżności
i warunków Cauchy'ego. Jeśli jedna norma jest zupełna,
każdy ciąg Cauchy'ego w drugiej ma granicę w pierwszej,
a zatem także w drugiej.
\end{proof}
\begin{twierdzenie}[Wymiar skończony]
\label{ma2:thm:finite}
W skończenie wymiarowej przestrzeni wszystkie normy są
równoważne i każda z nich jest zupełna.
\end{twierdzenie}
\begin{proof}
Dla przestrzeni zerowej tezy wynikają z istnienia wyłącznie
wektora zerowego. Dla wymiaru $d\ge1$ wybieramy bazę
$(e_j)_{j=1}^d$ i izomorfizm $Tz=\sum_j z_je_j$ z $\K^d$.
Dla dowolnej normy $p$ mamy $p(Tz)\le C\max_j|z_j|$,
gdzie $C=\sum_jp(e_j)>0$. Z odwrotnej nierówności trójkąta
funkcja $z\mapsto p(Tz)$ jest ciągła względem normy maksimum.
Na zwartej sferze $\max_j|z_j|=1$ (twierdzenie Heinego--Borela)
osiąga minimum $c>0$, bo $T$ jest różnowartościowe.
Skalowanie każdego niezerowego $z$ do tej sfery daje
$c\max_j|z_j|\le p(Tz)$; dla $z=0$ nierówność też zachodzi.
Porównując dwie dowolne normy z normą maksimum, dostajemy
ich równoważność. Jeśli $(Tz_n)$ jest ciągiem Cauchy'ego,
każda współrzędna $z_n$ jest ciągiem Cauchy'ego w $\K$.
Istnieje granica współrzędnościowa $z\in\K^d$; skończoność
liczby współrzędnych daje $\max_j|(z_n-z)_j|\to0$.
Górne oszacowanie daje $p(Tz_n-Tz)\to0$.
\end{proof}
\begin{przyklad}[Niezupełna przestrzeń unormowana]
Przestrzeń $c_{00}$ z normą supremum nie jest zupełna.
Ciąg $x^{(n)}=(1,1/2,\ldots,1/n,0,\ldots)$ jest ciągiem
Cauchy'ego, gdyż dla $m>n$ norma różnicy wynosi $1/(n+1)$.
W $\ell^\infty$ zbiega do $(1/j)_{j\ge1}$, ponieważ norma
różnicy od tego ciągu wynosi $1/(n+1)$. Jednoznaczność
granicy wyklucza granicę w $c_{00}$.
\end{przyklad}
